% ===============================================================
% Resume en francais de l'article
% « Teaching first-year engineering students to calibrate generative AI
%   with indirect prompt injections hidden in course materials »
% Selina Cheggour. (c) 2026, tous droits reserves.
% Compilation : pdflatex resume_fr.tex
% ===============================================================
\documentclass[11pt,a4paper]{article}
\usepackage[utf8]{inputenc}
\usepackage[T1]{fontenc}
\usepackage[french]{babel}
\usepackage{lmodern}
\usepackage{microtype}
\usepackage[a4paper,margin=2.5cm]{geometry}
\usepackage[hidelinks]{hyperref}

\hypersetup{
  pdftitle={Résumé en français : apprendre aux élèves ingénieurs de première année à étalonner l'IA générative avec des injections de prompt indirectes cachées dans les supports de cours},
  pdfauthor={Selina Cheggour},
  pdfkeywords={littératie en IA, formation d'ingénieurs, injection de prompt indirecte, IA générative, esprit critique, traçabilité, apprentissage par scénario}
}

\title{Apprendre aux élèves ingénieurs de première année à étalonner l'IA générative avec des
injections de prompt indirectes cachées dans les supports de cours}
\author{Selina Cheggour$^{1,2}$\\[4pt]
\small $^{1}$Université de Strasbourg, ICube UMR 7357, Strasbourg, France\\
\small $^{2}$Icam, Institut Catholique d'Arts et Métiers, Strasbourg, France\\[2pt]
\small \href{mailto:selina.cheggour@icam.fr}{selina.cheggour@icam.fr}}
\date{}

\begin{document}
\maketitle
\renewcommand{\thefootnote}{}
\footnotetext{Résumé en français de l'article « Teaching first-year engineering students to
calibrate generative AI with indirect prompt injections hidden in course materials ». Méthode
présentée au conseil de perfectionnement de l'Icam Toulouse les 14 et 15 octobre 2026.
\copyright{} 2026 Selina Cheggour. Tous droits réservés.}
\renewcommand{\thefootnote}{\arabic{footnote}}

\section*{Résumé}

Les élèves ingénieurs de première année utilisent déjà les assistants d'IA générative, et
l'aisance des réponses suscite une confiance qui n'a pas été méritée. Cet article présente une
séquence pédagogique qui traite un assistant d'IA comme un instrument de mesure à étalonner
avant de s'y fier. Une séance de travaux pratiques devient un atelier guidé, où les étudiants
soumettent des entrées connues à cinq outils d'IA et confrontent chaque sortie à une référence
qu'ils maîtrisent, suivi d'un mémo de révélation et d'un mini-projet décisionnel portés par un
scénario industriel continu. Les supports contiennent des injections de prompt indirectes,
cachées en texte de la couleur du fond, en micro-caractères dans une image et dans les
métadonnées du document, de sorte qu'un étudiant qui confie un document au modèle sans
l'inspecter est trompé comme le serait un professionnel. Le mini-projet est conçu pour qu'une
seule affirmation injectée inverse la bonne décision de sécurité, et l'évaluation repose sur la
traçabilité, avec une rubrique obligatoire de ce qui n'a pas pu être établi. Un essai à blanc
réalisé par l'enseignante avec un assistant grand public a montré que le piège peut agir comme
prévu, le modèle justifiant la mauvaise décision par des éléments absents de toute page
visible. La séquence offre aux formations d'ingénieurs une entrée sans prérequis dans la
sécurité de l'IA et la critique des sources.

\medskip
\noindent\textbf{Mots-clés :} littératie en IA, formation d'ingénieurs, injection de prompt
indirecte, IA générative, esprit critique, traçabilité, apprentissage par scénario

\section*{Pour citer l'article}

Cheggour, S. (2026). \emph{Teaching first-year engineering students to calibrate generative AI
with indirect prompt injections hidden in course materials}. Méthode présentée au conseil de
perfectionnement de l'Icam Toulouse, 14-15 octobre 2026. Dépôt de la méthodologie :
\href{https://doi.org/10.5281/zenodo.23390568}{doi.org/10.5281/zenodo.23390568}.

\end{document}
